\section{Introduction}\label{sec:intro}

\subsection{One idea, ten coats}

The operator-first transport programme that this paper belongs to is a series of preprints
\Cc{} around a single structure: a family of subspaces carried along a path, together with the
quantity that decides what can be resolved along the way. The series places that quantity in
different positions --- in the spectrum (Riesz projectors, exceptional points), between two
subspaces (principal angles, persistence), in the pseudospectrum (transient growth), at an
interface (refraction, Schur complements), or in the readout (numerator selection) --- and each
placement is a paper. What the series has not had is a statement of the structure that is
narrow enough to be wrong. This paper attempts one --- for three of its objects, and in a
different idiom: it uses the standard quantum geometric tensor and none of the programme's own
machinery (Dirac time, the vacuum-operator two-parameter family, connection energy). It is a
companion line, not a replacement for any part of the programme, and where it touches the
corpus's results it does so as an independent check.

\subsection{The thesis}

Two sentences, and the paper is their elaboration.

\begin{enumerate}[label=(\roman*)]
\item \textbf{Gaps.} A transport invariant is stable to the extent that the gap across
the boundary of the cluster defining it is open, and the relevant gap is never the one inside
the cluster.
\item \textbf{Totals.} A transport can forget almost every direction and still retain a constant;
for the distance map of $N$ bodies that constant is $N(N-2)$, and the shape only decides how it
is distributed.
\end{enumerate}

The first sentence is a correction of the series' ``hierarchy principle of structural
persistence'', which as stated is false at its first level. The second is a theorem the series
did not have, and it is a different statement from the per-shape spectral formula, not a
correction of it: that formula is about the map to squared distances and holds for its own
operator (Remark~\ref{rem:twomaps}).

\subsection{What is new here}

\begin{enumerate}[label=\textbf{C\arabic*.},leftmargin=2.2em]
\item \textbf{The redistribution theorem} (\S\ref{sec:redistribution}). The off-diagonal part of
a structural operator is traceless, vanishes if and only if the operator commutes with the
projector, and has the same Frobenius norm as the commutator. The algebra is elementary; the
point is that it makes ``redistribution'' and ``non-commutativity'' the same quantity, which the
series treats as two phenomena.

\item \textbf{The geometry is derived} (\S\ref{sec:geometry}). The quantum geometric tensor of a
projector family is $Q_{\mu\nu}=\Tr(\PP\partial_\mu\PP\partial_\nu\PP)$; its real part is a metric
because it is a squared norm of a block of $\dot\PP$, and the curvature is $-2\operatorname{Im}Q$
with $\abs{\omega}\le2\sqrt{g_{\mu\mu}g_{\nu\nu}}$. The object the series writes down,
$\Tr(\partial_\mu\PP\partial_\nu\PP)$, is real and is twice the metric; it cannot carry curvature.

\item \textbf{The hierarchy, quantified} (\S\ref{sec:hierarchy}). With constants: eigenvalues are
$1$-Lipschitz and unprotected; a rank-one projector moves at exactly the eigenvector's rate in
operator norm and $\sqrt2$ faster in Frobenius, so it protects nothing; a cluster projector is
protected by the external gap --- with the union gap in the denominator for a finite pair, and
constant $1$ in the infinitesimal form that the paper uses --- and its angle obeys a sharp
$\tan 2\Theta$ bound when the perturbation is off-diagonal and the cluster is a separated simple
eigenvalue; the $\pi/2$ that circulates in the
literature is the constant of the angle bound, attained at the edge of the stable range;
integers are locally constant. The exceptional point is invisible to the cluster and fatal to the basis.

\item \textbf{The capacity sum rule} (\S\ref{sec:totals}). $\sum_{i<j}\norm{\nabla d_{ij}}^2
=(N-1)\sum_i1/m_i-\sum_{i<j}d_{ij}^2$, hence $N(N-2)$ for equal unit masses, for every
configuration. The total is combinatorial; the distribution is the shape.

\item \textbf{The shape sphere, solved} (\S\ref{sec:shape}). $\sigma_1^2+\sigma_2^2=3$,
$\det M^*M=12A^2/(d_{12}^2d_{13}^2d_{23}^2)$ with maximum $9/4$ exactly at the equilateral shapes
and zero exactly on the collinear locus, a linear isotropic cone of slope $3$, rank drop where the
gap is largest, and holonomy $-1$. Mass symmetry breaking does not create the cone and does not
tilt it: it moves the two optic points off the equilateral shapes by $\varepsilon^2/9$ and makes
the cone elliptic.

\item \textbf{Refraction, demoted} (\S\ref{sec:refraction}). $\rho_1\sin^2\theta_1=\rho_2\sin^2
\theta_2$ is Snell's law for the conformal metric $\rho g$. The theorem is about geodesics; the
claim that a physical medium is conformal with factor $\rho$ is a hypothesis, and saying which is
which is the point.
\end{enumerate}

\subsection{What is not claimed}

No physical system is modelled here, and \S\ref{sec:scope} says so at the level of individual
claims: the covariance-data reproduction is a consistency check on published data and not
evidence for the framework; the cosmological reading is circular as stated and is not used; the
psychophysical reading is an interpretation without a theorem. The paper also does not claim that
the framework is new mathematics --- the algebra of \S\S\ref{sec:redistribution}--\ref{sec:geometry}
is elementary and the classical inputs of \S\ref{sec:hierarchy} are classical. What is new is the
assemblage, the constants of \S\ref{sec:hierarchy}, and the two exact results of
\S\S\ref{sec:totals}--\ref{sec:shape}. Several remarks record corrections; the reader should know
that some of them correct earlier drafts of \emph{this} paper rather than the corpus.

\subsection{How to read it}

Every statement carries a grade --- \T\ proved, \V\ machine-verified with the file and residual
named, \Cc\ cited, \Dd\ data, \Hh\ hypothesis --- and where a chain mixes grades the weakest link
is named in the sentence that draws the conclusion (\S\ref{sec:setting}). The paper is written to
be checked rather than believed: \S\ref{sec:refute} lists, for each claim, the cheapest
observation that would falsify it, and Appendix~\ref{app:repro} gives the commands that
re-run everything.

\subsection{Roadmap}

\S\ref{sec:setting} fixes notation and one symbol per concept. \S\ref{sec:redistribution} proves
the redistribution theorem, \S\ref{sec:geometry} derives the metric and curvature from it,
\S\ref{sec:hierarchy} states the hierarchy with its constants, crediting the corpus's own
tip-and-tail law, and records where earlier drafts of this paper were wrong.
\S\ref{sec:totals} proves the sum rule, \S\ref{sec:shape} solves the shape sphere,
\S\ref{sec:refraction} disposes of the refraction law. \S\S\ref{sec:scope}--\ref{sec:open} state
the scope, the refutations and the open problems; the appendices carry the ledger, the
reproduction instructions and the classical inputs.
